% Options for packages loaded elsewhere
\PassOptionsToPackage{unicode}{hyperref}
\PassOptionsToPackage{hyphens}{url}
\PassOptionsToPackage{dvipsnames,svgnames,x11names}{xcolor}
%
\documentclass[
  11pt,
  a4paper,
]{article}
\usepackage{amsmath,amssymb}
\usepackage{iftex}
\ifPDFTeX
  \usepackage[T1]{fontenc}
  \usepackage[utf8]{inputenc}
  \usepackage{textcomp} % provide euro and other symbols
\else % if luatex or xetex
  \usepackage{unicode-math} % this also loads fontspec
  \defaultfontfeatures{Scale=MatchLowercase}
  \defaultfontfeatures[\rmfamily]{Ligatures=TeX,Scale=1}
\fi
\usepackage{lmodern}
\ifPDFTeX\else
  % xetex/luatex font selection
\fi
% Use upquote if available, for straight quotes in verbatim environments
\IfFileExists{upquote.sty}{\usepackage{upquote}}{}
\IfFileExists{microtype.sty}{% use microtype if available
  \usepackage[]{microtype}
  \UseMicrotypeSet[protrusion]{basicmath} % disable protrusion for tt fonts
}{}
\makeatletter
\@ifundefined{KOMAClassName}{% if non-KOMA class
  \IfFileExists{parskip.sty}{%
    \usepackage{parskip}
  }{% else
    \setlength{\parindent}{0pt}
    \setlength{\parskip}{6pt plus 2pt minus 1pt}}
}{% if KOMA class
  \KOMAoptions{parskip=half}}
\makeatother
\usepackage{xcolor}
\usepackage{fancyhdr}
\usepackage{titlesec}
\usepackage{enumitem}
\usepackage[margin=25mm]{geometry}
\usepackage{longtable,booktabs,array}
\usepackage{calc} % for calculating minipage widths
% Correct order of tables after \paragraph or \subparagraph
\usepackage{etoolbox}
\makeatletter
\patchcmd\longtable{\par}{\if@noskipsec\mbox{}\fi\par}{}{}
\makeatother
% Allow footnotes in longtable head/foot
\IfFileExists{footnotehyper.sty}{\usepackage{footnotehyper}}{\usepackage{footnote}}
\makesavenoteenv{longtable}
\setlength{\emergencystretch}{3em} % prevent overfull lines
\providecommand{\tightlist}{%
  \setlength{\itemsep}{0pt}\setlength{\parskip}{0pt}}
\setcounter{secnumdepth}{3}
\setcounter{tocdepth}{2}
\ifLuaTeX
\usepackage[bidi=basic]{babel}
\else
\usepackage[bidi=default]{babel}
\fi
\babelprovide[main,import]{english}
% get rid of language-specific shorthands (see #6817):
\let\LanguageShortHands\languageshorthands
\def\languageshorthands#1{}
\ifLuaTeX
  \usepackage{selnolig}  % disable illegal ligatures
\fi
\usepackage{bookmark}
\IfFileExists{xurl.sty}{\usepackage{xurl}}{} % add URL line breaks if available
\urlstyle{same}
\hypersetup{
  pdftitle={Oriented Holomorphic Pure-Plus-Interaction Vector SRM},
  pdflang={en},
  colorlinks=true,
  linkcolor={MidnightBlue},
  filecolor={Maroon},
  citecolor={Blue},
  urlcolor={MidnightBlue},
  pdfcreator={LaTeX via pandoc}}

\title{Oriented Holomorphic Pure-Plus-Interaction Vector SRM}
\usepackage{etoolbox}
\makeatletter
\providecommand{\subtitle}[1]{% add subtitle to \maketitle
  \apptocmd{\@title}{\par {\large #1 \par}}{}{}
}
\makeatother
\subtitle{Standalone Technical Handoff}
\author{}
\date{}

\setlist[itemize]{leftmargin=*,itemsep=2pt,topsep=4pt}
\setlist[enumerate]{leftmargin=*,itemsep=2pt,topsep=4pt}
\titleformat{\section}{\Large\bfseries}{\thesection}{0.75em}{}
\titleformat{\subsection}{\large\bfseries}{\thesubsection}{0.75em}{}
\pagestyle{fancy}
\fancyhf{}
\fancyhead[L]{Oriented Holomorphic Vector SRM}
\fancyhead[R]{Technical Handoff}
\fancyfoot[C]{\thepage}
\setlength{\headheight}{14pt}
\setlength{\abovedisplayskip}{8pt plus 2pt minus 2pt}
\setlength{\belowdisplayskip}{8pt plus 2pt minus 2pt}
\begin{document}
\maketitle
\thispagestyle{empty}
\clearpage

{
\hypersetup{linkcolor=black}
\setcounter{tocdepth}{3}
\tableofcontents
}
\section{Objective}\label{objective}

Construct reduced complex Fourier coefficients

\[
\boldsymbol a(\boldsymbol k)\in\mathbb C^2
\]

for the two solenoidal polarization coordinates of every retained
Fourier mode, such that

\[
\mathbb E\!\left[
\boldsymbol a(\boldsymbol k)
\boldsymbol a(\boldsymbol p)^\dagger
\right]
=
\boldsymbol\Sigma(\boldsymbol k)\,
\delta_{\boldsymbol k\boldsymbol p},
\]

and, for every retained triad

\[
\boldsymbol q=\boldsymbol k+\boldsymbol p,
\]

\[
\mathbb E\!\left[
a_a(\boldsymbol k)
a_b(\boldsymbol p)
a_c(\boldsymbol q)^*
\right]
=
\mathcal B_{abc}(\boldsymbol k,\boldsymbol p).
\]

The generator must preserve the prescribed second-order covariance while
introducing exactly the required third-order phase coupling. It must
also remain valid when many triads share modes or feed the same output
mode.

The construction must reduce to an ordinary second-order solenoidal
spectral generator when

\[
\mathcal B=\boldsymbol0.
\]


\subsection{Fourier-space kinematics and reduced
coordinates}\label{fourier-space-kinematics-and-reduced-coordinates}

Let

\[
\widehat{\boldsymbol u}(\boldsymbol k)\in\mathbb C^3
\]

denote the Cartesian Fourier coefficient of the velocity field
\(\boldsymbol u\) at mode \(\boldsymbol k\). Let

\[
\boldsymbol d(\boldsymbol k)\in\mathbb C^3
\]

be the Fourier symbol of the selected discrete divergence operator,
defined so that

\begin{equation}
\widehat{(\nabla\cdot\boldsymbol u)}(\boldsymbol k)
=
\boldsymbol d(\boldsymbol k)^\dagger
\widehat{\boldsymbol u}(\boldsymbol k).
\label{eq:discrete-divergence-symbol}
\end{equation}

Accordingly, discrete solenoidality at mode \(\boldsymbol k\) is

\[
\boldsymbol d(\boldsymbol k)^\dagger
\widehat{\boldsymbol u}(\boldsymbol k)
=
0.
\]

Let

\[
\boldsymbol Q(\boldsymbol k)\in\mathbb C^{3\times2}
\]

have orthonormal columns spanning the transverse subspace

\[
\operatorname{range}\boldsymbol Q(\boldsymbol k)
=
\ker\!\left(\boldsymbol d(\boldsymbol k)^\dagger\right).
\]

Thus,

\[
\boldsymbol Q(\boldsymbol k)^\dagger\boldsymbol Q(\boldsymbol k)
=
\boldsymbol I_2,
\qquad
\boldsymbol d(\boldsymbol k)^\dagger\boldsymbol Q(\boldsymbol k)
=
\boldsymbol0.
\]

The reduced polarization coefficient \(\boldsymbol a(\boldsymbol k)\) and
the Cartesian Fourier coefficient are related by

\[
\widehat{\boldsymbol u}(\boldsymbol k)
=
\boldsymbol Q(\boldsymbol k)\boldsymbol a(\boldsymbol k).
\]

\section{Main difficulty}\label{main-difficulty}

A quadratic interaction consumes covariance power and may generate
unintended contractions.

For a generic split

\[
\boldsymbol a(\boldsymbol q)
=
\boldsymbol a^{(p)}(\boldsymbol q)
+
\boldsymbol a^{(I)}(\boldsymbol q),
\]

the complete covariance is

\[
\boldsymbol\Sigma
=
\boldsymbol\Sigma^{(p)}
+
\boldsymbol\Sigma^{(I)}
+
\boldsymbol\Sigma^{(pI)}
+
\boldsymbol\Sigma^{(Ip)},
\]

and the complete third moment contains eight combinations of pure and
interaction terms.

A valid architecture must therefore ensure that:

\begin{enumerate}
\def\labelenumi{\arabic{enumi}.}
\tightlist
\item
  the intended pure--pure--interaction term produces the target
  bispectrum;
\item
  the pure--interaction covariance terms vanish;
\item
  interaction covariance includes all parent-pair cross contractions;
\item
  unintended third-order terms vanish, including
  interaction--interaction--interaction terms;
\item
  the residual pure covariance remains positive semidefinite;
\item
  overlapping triads do not create uncontrolled moment contributions.
\end{enumerate}

\section{Independent half-spectrum and triad
orientation}\label{independent-half-spectrum-and-triad-orientation}

Let \(\mathcal H\) contain one representative from each
non-self-conjugate Fourier pair

\[
\{\boldsymbol k,-\boldsymbol k\}.
\]

Choose a positive half-spectrum and a strict order \(\prec\) such that,
for every retained non-aliased relation

\[
\boldsymbol q=\boldsymbol k+\boldsymbol p\in\mathcal H,
\]

one has

\[
\boldsymbol k\prec\boldsymbol q,
\qquad
\boldsymbol p\prec\boldsymbol q.
\]

For each output mode \(\boldsymbol q\), define the set of canonical
parent pairs

\[
\mathcal P_{\boldsymbol q}
=
\left\{
(\boldsymbol k,\boldsymbol p):
\boldsymbol k+\boldsymbol p=\boldsymbol q,\;
\boldsymbol k,\boldsymbol p\in\mathcal H,\;
\boldsymbol k\preceq\boldsymbol p
\right\}.
\]

Only one oriented output leg is assigned to each canonical triad: the
interaction generated by \((\boldsymbol k,\boldsymbol p)\) is inserted
into the child mode \(\boldsymbol q=\boldsymbol k+\boldsymbol p\).

This orientation makes the construction triangular in Fourier order and
allows all interaction terms to be holomorphic functions of the
independent complex Gaussian variables.

\section{Gaussian carrier variables}\label{gaussian-carrier-variables}

For each \(\boldsymbol k\in\mathcal H\), introduce an independent proper
complex Gaussian vector

\[
\boldsymbol z(\boldsymbol k)\sim
\mathcal{CN}(\boldsymbol0,\boldsymbol I_2).
\]

The carrier variables satisfy

\[
\mathbb E[\boldsymbol z(\boldsymbol k)]
=
\boldsymbol0,
\]

\[
\mathbb E\!\left[
\boldsymbol z(\boldsymbol k)
\boldsymbol z(\boldsymbol p)^\dagger
\right]
=
\delta_{\boldsymbol k\boldsymbol p}\boldsymbol I_2,
\]

and

\[
\mathbb E\!\left[
\boldsymbol z(\boldsymbol k)
\boldsymbol z(\boldsymbol p)^T
\right]
=
\boldsymbol0.
\]

Variables associated with different half-spectrum modes are independent.

The same carrier family must drive both the pure and interaction terms.
This shared-carrier structure is necessary for a nonzero
pure--pure--interaction third moment.

\section{Pure term}\label{pure-term}

For every mode,

\[
\boldsymbol a^{(p)}(\boldsymbol k)
=
\boldsymbol L(\boldsymbol k)
\boldsymbol z(\boldsymbol k),
\]

where

\[
\boldsymbol L(\boldsymbol k)
\boldsymbol L(\boldsymbol k)^\dagger
=
\boldsymbol\Sigma^{(p)}(\boldsymbol k).
\]

The pure covariance \(\boldsymbol\Sigma^{(p)}\) is not prescribed
independently. It is the residual covariance remaining after the
interaction covariance has been deducted from the target covariance.

\section{Quadratic interaction term}\label{quadratic-interaction-term}

\subsection{1 Distinct parent modes}\label{distinct-parent-modes}

For a canonical pair \(\boldsymbol k\ne\boldsymbol p\), define the
quadratic feature

\[
\boldsymbol\phi_{\boldsymbol k\boldsymbol p}
=
\boldsymbol z(\boldsymbol k)
\otimes
\boldsymbol z(\boldsymbol p)
\in\mathbb C^4.
\]

The interaction contribution to output mode \(\boldsymbol q\) is

\[
\boldsymbol a^{(I)}(\boldsymbol q)
=
\sum_{(\boldsymbol k,\boldsymbol p)\in\mathcal P_{\boldsymbol q}}
\boldsymbol A(\boldsymbol q;\boldsymbol k,\boldsymbol p)
\boldsymbol\phi_{\boldsymbol k\boldsymbol p},
\]

where

\[
\boldsymbol A(\boldsymbol q;\boldsymbol k,\boldsymbol p)
\in\mathbb C^{2\times4}.
\]

Equivalently,

\[
a_c^{(I)}(\boldsymbol q)
=
\sum_{(\boldsymbol k,\boldsymbol p)\in\mathcal P_{\boldsymbol q}}
\sum_{i,j=1}^{2}
A_{cij}(\boldsymbol q;\boldsymbol k,\boldsymbol p)
z_i(\boldsymbol k)z_j(\boldsymbol p).
\]

Because the carriers are centered proper complex Gaussians,

\[
\mathbb E[
z_i(\boldsymbol k)z_j(\boldsymbol p)
]
=
0.
\]

Therefore every interaction term is centered.

\subsection{2 Repeated parent mode}\label{repeated-parent-mode}

For a repeated-parent triad

\[
\boldsymbol q=2\boldsymbol k,
\]

use the normalized symmetric second-order feature

\[
\boldsymbol h_2(\boldsymbol z)
=
\begin{bmatrix}
z_1^2/\sqrt2\\
z_1z_2\\
z_2^2/\sqrt2
\end{bmatrix}.
\]

It satisfies

\[
\mathbb E[
\boldsymbol h_2(\boldsymbol z)
\boldsymbol h_2(\boldsymbol z)^\dagger
]
=
\boldsymbol I_3.
\]

The coefficient matrix for a repeated-parent interaction therefore
belongs to

\[
\mathbb C^{2\times3}.
\]

The target bispectrum must satisfy the parent-exchange symmetry

\[
\mathcal B_{abc}(\boldsymbol k,\boldsymbol k)
=
\mathcal B_{bac}(\boldsymbol k,\boldsymbol k).
\]

\section{Complete second-order
moment}\label{complete-second-order-moment}

For any output mode,

\[
\boldsymbol a(\boldsymbol q)
=
\boldsymbol a^{(p)}(\boldsymbol q)
+
\boldsymbol a^{(I)}(\boldsymbol q).
\]

Hence

\[
\begin{aligned}
\mathbb E[
\boldsymbol a(\boldsymbol q)
\boldsymbol a(\boldsymbol q)^\dagger
]
={}&
\boldsymbol\Sigma^{(p)}(\boldsymbol q)
+
\boldsymbol\Sigma^{(I)}(\boldsymbol q)
\\
&+
\boldsymbol\Sigma^{(pI)}(\boldsymbol q)
+
\boldsymbol\Sigma^{(Ip)}(\boldsymbol q).
\end{aligned}
\]

\subsection{1 Vanishing pure--interaction
covariance}\label{vanishing-pureinteraction-covariance}

The cross term contains a third-order moment of a centered proper
complex Gaussian family:

\[
\mathbb E\!\left[
\boldsymbol z(\boldsymbol q)
\boldsymbol\phi_{\boldsymbol k\boldsymbol p}^\dagger
\right]
=
\boldsymbol0.
\]

Therefore,

\[
\boldsymbol\Sigma^{(pI)}(\boldsymbol q)
=
\boldsymbol0,
\qquad
\boldsymbol\Sigma^{(Ip)}(\boldsymbol q)
=
\boldsymbol0.
\]

Thus

\[
\boldsymbol\Sigma(\boldsymbol q)
=
\boldsymbol\Sigma^{(p)}(\boldsymbol q)
+
\boldsymbol\Sigma^{(I)}(\boldsymbol q).
\]

\subsection{2 General interaction
covariance}\label{general-interaction-covariance}

For ordered distinct parent pairs,

\[
\begin{aligned}
&
\mathbb E\!\left[
z_i(\boldsymbol k)z_j(\boldsymbol p)
z_m(\boldsymbol r)^*z_n(\boldsymbol s)^*
\right]
\\
&\quad=
\delta_{\boldsymbol k\boldsymbol r}
\delta_{\boldsymbol p\boldsymbol s}
\delta_{im}\delta_{jn}
+
\delta_{\boldsymbol k\boldsymbol s}
\delta_{\boldsymbol p\boldsymbol r}
\delta_{in}\delta_{jm}.
\end{aligned}
\]

Let \(\alpha,\beta\in\mathcal P_{\boldsymbol q}\) denote canonical
parent pairs and define

\[
\boldsymbol\Gamma_{\alpha\beta}
=
\mathbb E[
\boldsymbol\phi_\alpha
\boldsymbol\phi_\beta^\dagger
].
\]

Then the complete interaction covariance is

\[
\boldsymbol\Sigma^{(I)}(\boldsymbol q)
=
\sum_{\alpha,\beta\in\mathcal P_{\boldsymbol q}}
\boldsymbol A_{\boldsymbol q,\alpha}
\boldsymbol\Gamma_{\alpha\beta}
\boldsymbol A_{\boldsymbol q,\beta}^\dagger.
\]

This expression includes every cross contraction between parent pairs
feeding the same output mode.

After canonicalizing unordered pairs and using normalized
repeated-parent features,

\[
\boldsymbol\Gamma_{\alpha\beta}
=
\delta_{\alpha\beta}\boldsymbol I.
\]

Therefore,

\[
\boxed{
\boldsymbol\Sigma^{(I)}(\boldsymbol q)
=
\sum_{\alpha\in\mathcal P_{\boldsymbol q}}
\boldsymbol A_{\boldsymbol q,\alpha}
\boldsymbol A_{\boldsymbol q,\alpha}^\dagger
}.
\]

Different parent pairs may share one Fourier mode, but their quadratic
features remain orthogonal unless their complete parent multisets
coincide.

\subsection{3 Residual covariance
condition}\label{residual-covariance-condition}

The pure covariance is

\[
\boxed{
\boldsymbol\Sigma^{(p)}(\boldsymbol q)
=
\boldsymbol\Sigma(\boldsymbol q)
-
\boldsymbol\Sigma^{(I)}(\boldsymbol q)
}.
\]

The construction is admissible only if

\[
\boxed{
\boldsymbol\Sigma^{(p)}(\boldsymbol q)
\succeq\boldsymbol0
}
\]

for every retained output mode.

The interaction therefore consumes covariance power, and the prescribed
covariance and bispectrum cannot be chosen independently.

\section{Complete third-order moment}\label{complete-third-order-moment}

For a retained triad

\[
\boldsymbol q=\boldsymbol k+\boldsymbol p,
\]

expand

\[
\mathbb E\!\left[
(a_a^{(p)}+a_a^{(I)})_{\boldsymbol k}
(a_b^{(p)}+a_b^{(I)})_{\boldsymbol p}
(a_c^{(p)}+a_c^{(I)})_{\boldsymbol q}^*
\right].
\]

There are eight contributions:

\[
PPP^*,\;
PPI^*,\;
PIP^*,\;
IPP^*,\;
PII^*,\;
IPI^*,\;
IIP^*,\;
III^*.
\]

Under the global phase transformation

\[
\boldsymbol z(\boldsymbol r)
\longmapsto
e^{\mathrm i\theta}
\boldsymbol z(\boldsymbol r)
\]

for every half-spectrum mode \(\boldsymbol r\), the joint carrier
distribution is unchanged.

The pure and interaction terms transform as

\[
\boldsymbol a^{(p)}
\longmapsto
e^{\mathrm i\theta}\boldsymbol a^{(p)},
\]

\[
\boldsymbol a^{(I)}
\longmapsto
e^{2\mathrm i\theta}\boldsymbol a^{(I)}.
\]

The phase exponents of the eight third-moment terms are:

\begin{longtable}[]{@{}lr@{}}
\toprule\noalign{}
Term & Phase exponent \\
\midrule\noalign{}
\endhead
\bottomrule\noalign{}
\endlastfoot
\(PPP^*\) & \(1\) \\
\(PPI^*\) & \(0\) \\
\(PIP^*\) & \(2\) \\
\(IPP^*\) & \(2\) \\
\(PII^*\) & \(1\) \\
\(IPI^*\) & \(1\) \\
\(IIP^*\) & \(3\) \\
\(III^*\) & \(2\) \\
\end{longtable}

Only \(PPI^*\) is invariant for every \(\theta\). Therefore,

\[
\boxed{
\mathbb E[
a_a(\boldsymbol k)
a_b(\boldsymbol p)
a_c(\boldsymbol q)^*
]
=
\mathbb E[
a_a^{(p)}(\boldsymbol k)
a_b^{(p)}(\boldsymbol p)
a_c^{(I)}(\boldsymbol q)^*
]
}.
\]

All unintended third-order contributions vanish identically, including

\[
\mathbb E[
a_a^{(I)}(\boldsymbol k)
a_b^{(I)}(\boldsymbol p)
a_c^{(I)}(\boldsymbol q)^*
]
=
0.
\]

This remains true even when many triads overlap, because the argument
depends on holomorphic degree rather than pairwise triad independence.

\section{Relation between the interaction tensor and the target
bispectrum}\label{relation-between-the-interaction-tensor-and-the-target-bispectrum}

For a distinct parent pair, define the output-component matrix

\[
\boldsymbol A_c(\boldsymbol q;\boldsymbol k,\boldsymbol p)
=
\left[
A_{cij}(\boldsymbol q;\boldsymbol k,\boldsymbol p)
\right]_{i,j=1}^{2}
\in\mathbb C^{2\times2}.
\]

Define the corresponding target bispectrum slice

\[
\boldsymbol B_c(\boldsymbol k,\boldsymbol p)
=
\left[
\mathcal B_{abc}(\boldsymbol k,\boldsymbol p)
\right]_{a,b=1}^{2}.
\]

Wick contraction gives

\[
\boxed{
\boldsymbol B_c(\boldsymbol k,\boldsymbol p)
=
\boldsymbol L(\boldsymbol k)\,
\overline{
\boldsymbol A_c(\boldsymbol q;\boldsymbol k,\boldsymbol p)
}
\,\boldsymbol L(\boldsymbol p)^T
}.
\]

If the pure covariance factors are nonsingular,

\[
\boxed{
\boldsymbol A_c(\boldsymbol q;\boldsymbol k,\boldsymbol p)
=
\overline{
\boldsymbol L(\boldsymbol k)^{-1}
\boldsymbol B_c(\boldsymbol k,\boldsymbol p)
\boldsymbol L(\boldsymbol p)^{-T}
}
}.
\]

For rank-deficient pure covariances, use Moore--Penrose inverses:

\[
\boxed{
\boldsymbol A_c
=
\overline{
\boldsymbol L(\boldsymbol k)^+
\boldsymbol B_c
\boldsymbol L(\boldsymbol p)^{+T}
}
}.
\]

The target slice must lie in the supported parent subspaces:

\[
\boxed{
\boldsymbol B_c
=
\boldsymbol P_{\boldsymbol k}
\boldsymbol B_c
\boldsymbol P_{\boldsymbol p}^T
},
\]

where

\[
\boldsymbol P_{\boldsymbol k}
=
\boldsymbol L(\boldsymbol k)
\boldsymbol L(\boldsymbol k)^+.
\]

Any interaction component lying in a null direction of the moment map
should be set to zero because it consumes covariance without changing
the prescribed bispectrum.

\section{Covariance cost of a target bispectrum
slice}\label{covariance-cost-of-a-target-bispectrum-slice}

For a distinct parent pair, let

\[
\boldsymbol G_{\boldsymbol q;\boldsymbol k\boldsymbol p}
=
\boldsymbol A_{\boldsymbol q;\boldsymbol k\boldsymbol p}
\boldsymbol A_{\boldsymbol q;\boldsymbol k\boldsymbol p}^\dagger.
\]

Its entries can be expressed directly in terms of the target bispectrum
and the parent pure covariances:

\[
\boxed{
\left[
\boldsymbol G_{\boldsymbol q;\boldsymbol k\boldsymbol p}
\right]_{cd}
=
\operatorname{tr}\!\left[
\boldsymbol B_c^\dagger
\boldsymbol\Sigma^{(p)}(\boldsymbol k)^+
\boldsymbol B_d
\boldsymbol\Sigma^{(p)}(\boldsymbol p)^{+T}
\right]
}.
\]

Hence

\[
\boxed{
\boldsymbol\Sigma^{(p)}(\boldsymbol q)
=
\boldsymbol\Sigma(\boldsymbol q)
-
\sum_{\boldsymbol k+\boldsymbol p=\boldsymbol q}
\boldsymbol G_{\boldsymbol q;\boldsymbol k\boldsymbol p}
}.
\]

This formula makes the covariance--bispectrum realizability coupling
explicit.

\section{Direct recursive
construction}\label{direct-recursive-construction}

The oriented half-spectrum architecture yields a triangular recursion.

Process modes in increasing order under \(\prec\).

For each output mode \(\boldsymbol q\):

\begin{enumerate}
\def\labelenumi{\arabic{enumi}.}
\item
  Enumerate all canonical parent pairs

  \[
  (\boldsymbol k,\boldsymbol p)\in\mathcal P_{\boldsymbol q}.
  \]
\item
  The parent pure factors

  \[
  \boldsymbol L(\boldsymbol k),
  \qquad
  \boldsymbol L(\boldsymbol p)
  \]

  are already available because

  \[
  \boldsymbol k,\boldsymbol p\prec\boldsymbol q.
  \]
\item
  Compute every interaction tensor

  \[
  \boldsymbol A(\boldsymbol q;\boldsymbol k,\boldsymbol p)
  \]

  from the prescribed target slice \(\boldsymbol B_c\).
\item
  Accumulate the interaction covariance:

  \[
  \boldsymbol\Sigma^{(I)}(\boldsymbol q)
  =
  \sum_{(\boldsymbol k,\boldsymbol p)\in\mathcal P_{\boldsymbol q}}
  \boldsymbol A_{\boldsymbol q;\boldsymbol k\boldsymbol p}
  \boldsymbol A_{\boldsymbol q;\boldsymbol k\boldsymbol p}^\dagger.
  \]
\item
  Form the residual:

  \[
  \boldsymbol\Sigma^{(p)}(\boldsymbol q)
  =
  \boldsymbol\Sigma(\boldsymbol q)
  -
  \boldsymbol\Sigma^{(I)}(\boldsymbol q).
  \]
\item
  Check the realizability condition:

  \[
  \lambda_{\min}
  \bigl(
  \boldsymbol\Sigma^{(p)}(\boldsymbol q)
  \bigr)
  \ge
  -\varepsilon_{\mathrm{PSD}}.
  \]
\item
  Reject the target if the negative eigenvalue is materially larger than
  numerical tolerance.
\item
  Clip only roundoff-scale negative eigenvalues.
\item
  Compute a square root:

  \[
  \boldsymbol\Sigma^{(p)}(\boldsymbol q)
  =
  \boldsymbol L(\boldsymbol q)
  \boldsymbol L(\boldsymbol q)^\dagger.
  \]
\end{enumerate}

No global nonlinear iteration is required for the non-aliased oriented
architecture.

A coupled solve becomes necessary only when introducing modular
aliasing, symmetric interaction allocation over multiple legs,
nonholomorphic interactions, or exceptional self-conjugate modes.

\section{Sample generator}\label{sample-generator}

After the offline recursive construction has determined all
\(\boldsymbol L\) and \(\boldsymbol A\):

\begin{enumerate}
\def\labelenumi{\arabic{enumi}.}
\item
  Draw independent carriers

  \[
  \boldsymbol z(\boldsymbol k)
  \sim\mathcal{CN}(\boldsymbol0,\boldsymbol I_2),
  \qquad
  \boldsymbol k\in\mathcal H.
  \]
\item
  Compute the pure terms

  \[
  \boldsymbol a^{(p)}(\boldsymbol k)
  =
  \boldsymbol L(\boldsymbol k)\boldsymbol z(\boldsymbol k).
  \]
\item
  Process modes in increasing order and compute

  \[
  \boldsymbol a^{(I)}(\boldsymbol q)
  =
  \sum_{(\boldsymbol k,\boldsymbol p)\in\mathcal P_{\boldsymbol q}}
  \boldsymbol A(\boldsymbol q;\boldsymbol k,\boldsymbol p)
  \boldsymbol\phi_{\boldsymbol k\boldsymbol p}.
  \]
\item
  Set

  \[
  \boldsymbol a(\boldsymbol q)
  =
  \boldsymbol a^{(p)}(\boldsymbol q)
  +
  \boldsymbol a^{(I)}(\boldsymbol q).
  \]
\item
  Reconstruct the negative half-spectrum through Fourier reality.
\end{enumerate}

The interaction values do not need to be used recursively as carriers.
Every interaction term is evaluated directly from the underlying
Gaussian variables \(\boldsymbol z\).

\section{Fourier reality and Cartesian
reconstruction}\label{fourier-reality-and-cartesian-reconstruction}

Using the Fourier-space notation introduced in
Section~\ref{fourier-space-kinematics-and-reduced-coordinates}, reconstruct
the Cartesian Fourier coefficient by

\[
\widehat{\boldsymbol u}(\boldsymbol k)
=
\boldsymbol Q(\boldsymbol k)
\boldsymbol a(\boldsymbol k),
\]

where the columns of \(\boldsymbol Q(\boldsymbol k)\) form an orthonormal
basis of

\[
\ker\!\left(\boldsymbol d(\boldsymbol k)^\dagger\right).
\]

If the basis gauge is chosen so that

\[
\boldsymbol Q(-\boldsymbol k)
=
\overline{\boldsymbol Q(\boldsymbol k)},
\]

impose

\[
\boldsymbol a(-\boldsymbol k)
=
\overline{\boldsymbol a(\boldsymbol k)}.
\]

Then

\[
\widehat{\boldsymbol u}(-\boldsymbol k)
=
\overline{\widehat{\boldsymbol u}(\boldsymbol k)}.
\]

For a general transverse-basis gauge, define

\[
\boldsymbol J(\boldsymbol k)
=
\boldsymbol Q(-\boldsymbol k)^\dagger
\overline{\boldsymbol Q(\boldsymbol k)}
\]

and impose

\[
\boldsymbol a(-\boldsymbol k)
=
\boldsymbol J(\boldsymbol k)
\overline{\boldsymbol a(\boldsymbol k)}.
\]

The Cartesian covariance is

\[
\boldsymbol\Phi(\boldsymbol k)
=
\boldsymbol Q(\boldsymbol k)
\boldsymbol\Sigma(\boldsymbol k)
\boldsymbol Q(\boldsymbol k)^\dagger.
\]

The Cartesian bispectrum is

\[
\mathbb E[
\widehat u_i(\boldsymbol k)
\widehat u_j(\boldsymbol p)
\widehat u_\ell(\boldsymbol q)^*
]
=
\sum_{a,b,c=1}^{2}
Q_{ia}(\boldsymbol k)
Q_{jb}(\boldsymbol p)
Q_{\ell c}(\boldsymbol q)^*
\mathcal B_{abc}(\boldsymbol k,\boldsymbol p).
\]

Solenoidality follows directly from the transverse basis.

\section{Exceptional cases}\label{exceptional-cases}

\subsection{1 Zero mode}\label{zero-mode}

The zero mode should normally represent the deterministic mean and be
excluded from the fluctuating generator:

\[
\boldsymbol a(\boldsymbol0)
=
\boldsymbol0.
\]

Triads containing the zero mode should not be included in the core
interaction set.

\subsection{2 Self-conjugate and Nyquist
modes}\label{self-conjugate-and-nyquist-modes}

On an even discrete grid, some modes satisfy

\[
-\boldsymbol k\equiv\boldsymbol k.
\]

Their Fourier coefficients are real rather than proper complex Gaussian.
The global complex-phase selection rule does not apply to them.

The preferred baseline is to use an odd or suitably dealiased grid that
excludes such modes from the interacting set.

A conservative alternative is to generate self-conjugate modes linearly
and prescribe zero bispectrum for triads containing them.

If such triads must be retained, use centered real Hermite features and
explicitly retain all Wick contractions in the governing equations.

\subsection{3 Repeated parent modes}\label{repeated-parent-modes}

For

\[
\boldsymbol q=2\boldsymbol k,
\]

use the normalized symmetric quadratic feature described above. Solve
the moment equation on the symmetric tensor-product subspace.

\subsection{4 Aliasing}\label{aliasing}

The triangular recursion assumes ordinary wavevector addition:

\[
\boldsymbol q=\boldsymbol k+\boldsymbol p.
\]

It does not apply directly to modular alias relations.

Aliased interactions should either be removed by dealiasing or handled
by a global coupled system.

\subsection{5 Rank-deficient parent
covariance}\label{rank-deficient-parent-covariance}

When

\[
\operatorname{rank}
\boldsymbol\Sigma^{(p)}(\boldsymbol k)<2,
\]

the target bispectrum must lie in the supported parent subspace. Use
pseudoinverses and explicitly test the support projector condition.

\section{Architecture rationale}\label{architecture-rationale}

The selected architecture has four central advantages.

\subsection{Exact third-moment
selection}\label{exact-third-moment-selection}

The global phase-degree argument eliminates all unwanted third-order
contributions exactly.

\subsection{Controlled overlap}\label{controlled-overlap}

Overlapping triads do not produce unintended third moments, and
canonical quadratic features diagonalize the interaction covariance
across distinct parent pairs.

\subsection{Low latent dimension}\label{low-latent-dimension}

Only two proper complex Gaussian carrier variables are needed per
independent Fourier mode.

\subsection{Direct solution}\label{direct-solution}

The half-spectrum ordering turns the covariance--bispectrum coupling
into a recursive construction instead of a global nonlinear moment
problem.

The main limitation is that the stochastic generator is explicitly
oriented: the interaction is assigned to one output leg of each
canonical triad. The resulting covariance and bispectrum can still
satisfy the required physical permutation and Fourier-reality
symmetries, but those symmetries must be imposed on the target data and
on the canonical triad bookkeeping.

\section{Implementation structure}\label{implementation-structure}

A practical implementation should separate the following components.

\subsection{1 Spectral geometry}\label{spectral-geometry}

\begin{itemize}
\tightlist
\item
  construct the independent half-spectrum;
\item
  identify zero, Nyquist, and self-conjugate modes;
\item
  impose a strict parent-before-child ordering;
\item
  enumerate canonical parent pairs;
\item
  reject or flag aliased triads.
\end{itemize}

\subsection{2 Target validation}\label{target-validation}

\begin{itemize}
\tightlist
\item
  check Hermitian positive semidefiniteness of every
  \(\boldsymbol\Sigma(\boldsymbol k)\);
\item
  check bispectrum parent-exchange symmetry;
\item
  check Fourier-reality relations;
\item
  check support consistency for rank-deficient parent covariances;
\item
  check repeated-parent symmetry.
\end{itemize}

\subsection{3 Offline coefficient
construction}\label{offline-coefficient-construction}

\begin{itemize}
\tightlist
\item
  recursively compute \(\boldsymbol A\);
\item
  accumulate \(\boldsymbol\Sigma^{(I)}\);
\item
  compute and factor \(\boldsymbol\Sigma^{(p)}\);
\item
  record PSD margins and condition numbers.
\end{itemize}

\subsection{4 Sampling}\label{sampling}

\begin{itemize}
\tightlist
\item
  draw independent proper complex Gaussian carriers;
\item
  evaluate pure coefficients;
\item
  evaluate quadratic interactions;
\item
  reconstruct the full Fourier spectrum;
\item
  transform to Cartesian coefficients or physical space.
\end{itemize}

\subsection{5 Analytical diagnostics}\label{analytical-diagnostics}

\begin{itemize}
\tightlist
\item
  compute generated covariance directly from \(\boldsymbol L\) and
  \(\boldsymbol A\);
\item
  compute generated bispectrum directly from the Wick-contraction
  formulas;
\item
  compare analytical moments against prescribed targets before Monte
  Carlo sampling.
\end{itemize}

\section{Required unit tests}\label{required-unit-tests}

The implementation should include at least the following tests.

\begin{enumerate}
\def\labelenumi{\arabic{enumi}.}
\item
  \textbf{Carrier statistics}

  Verify

  \[
  \mathbb E[
  \boldsymbol z(\boldsymbol k)
  \boldsymbol z(\boldsymbol p)^\dagger
  ]
  =
  \delta_{\boldsymbol k\boldsymbol p}\boldsymbol I_2
  \]

  and

  \[
  \mathbb E[
  \boldsymbol z(\boldsymbol k)
  \boldsymbol z(\boldsymbol p)^T
  ]
  =
  \boldsymbol0.
  \]
\item
  \textbf{Fourier reality}

  Verify reconstructed negative modes satisfy the required conjugacy
  relation.
\item
  \textbf{Pure--interaction covariance}

  Verify analytically and numerically that

  \[
  \boldsymbol\Sigma^{(pI)}
  =
  \boldsymbol\Sigma^{(Ip)}
  =
  \boldsymbol0.
  \]
\item
  \textbf{Quadratic-feature Gram matrix}

  Verify orthonormality for distinct canonical pairs and repeated-parent
  features.
\item
  \textbf{Multiple parent pairs for one output}

  Use a manufactured mode \(\boldsymbol q\) with several parent
  decompositions and verify additive interaction covariance.
\item
  \textbf{Overlapping triads}

  Use several triads sharing one or more modes and verify that no
  unintended third-order terms appear.
\item
  \textbf{Repeated-parent triad}

  Test \(\boldsymbol q=2\boldsymbol k\) using the symmetric quadratic
  basis.
\item
  \textbf{Rank-deficient covariance}

  Verify support projection, pseudoinverse handling, and rejection of
  unsupported bispectrum components.
\item
  \textbf{Residual PSD failure}

  Verify that materially indefinite residual covariance is detected and
  rejected.
\item
  \textbf{Roundoff-scale eigenvalues}

  Verify that only numerical-scale negative eigenvalues are clipped.
\item
  \textbf{Cartesian solenoidality}

  Verify

  \[
  \boldsymbol d(\boldsymbol k)^\dagger
  \widehat{\boldsymbol u}(\boldsymbol k)
  =
  0.
  \]
\item
  \textbf{Zero-bispectrum limit}

  Verify

  \[
  \mathcal B=\boldsymbol0
  \]

  implies

  \[
  \boldsymbol A=\boldsymbol0,
  \qquad
  \boldsymbol\Sigma^{(p)}=\boldsymbol\Sigma,
  \qquad
  \boldsymbol a=\boldsymbol L\boldsymbol z.
  \]
\end{enumerate}

\section{Monte Carlo validation}\label{monte-carlo-validation}

For manufactured targets known to be realizable by the architecture:

\begin{enumerate}
\def\labelenumi{\arabic{enumi}.}
\item
  construct \(\boldsymbol L\) and \(\boldsymbol A\) analytically;
\item
  generate a large ensemble of independent samples;
\item
  estimate

  \[
  \widehat{\boldsymbol\Sigma}(\boldsymbol k)
  =
  \frac1N
  \sum_{n=1}^{N}
  \boldsymbol a^{(n)}(\boldsymbol k)
  \boldsymbol a^{(n)}(\boldsymbol k)^\dagger;
  \]
\item
  estimate

  \[
  \widehat{\mathcal B}_{abc}(\boldsymbol k,\boldsymbol p)
  =
  \frac1N
  \sum_{n=1}^{N}
  a_a^{(n)}(\boldsymbol k)
  a_b^{(n)}(\boldsymbol p)
  a_c^{(n)}(\boldsymbol q)^*;
  \]
\item
  compare the empirical moments against both the prescribed targets and
  the analytical generated moments;
\item
  assess errors relative to estimated Monte Carlo standard errors rather
  than using only fixed absolute tolerances.
\end{enumerate}

The analytical formulas should match the prescribed moments to numerical
precision. Monte Carlo estimates should converge within sampling
tolerances.

\section{Completion criterion}\label{completion-criterion}

The construction is complete when, for every manufactured target known
to be realizable:

\[
\boldsymbol\Sigma_{\mathrm{generated}}(\boldsymbol k)
=
\boldsymbol\Sigma_{\mathrm{target}}(\boldsymbol k)
\]

and

\[
\mathcal B_{\mathrm{generated}}(\boldsymbol k,\boldsymbol p)
=
\mathcal B_{\mathrm{target}}(\boldsymbol k,\boldsymbol p)
\]

hold analytically to numerical precision, while Monte Carlo estimates
recover both moments within sampling uncertainty.

In the zero-bispectrum limit,

\[
\mathcal B=\boldsymbol0,
\]

the construction must reduce exactly to

\[
\boldsymbol a(\boldsymbol k)
=
\boldsymbol L(\boldsymbol k)\boldsymbol z(\boldsymbol k),
\]

with

\[
\boldsymbol L(\boldsymbol k)
\boldsymbol L(\boldsymbol k)^\dagger
=
\boldsymbol\Sigma(\boldsymbol k).
\]

\section{Recommended next implementation
step}\label{recommended-next-implementation-step}

Implement the non-aliased, non-self-conjugate, distinct-parent case
first.

The first manufactured validation problem should contain:

\begin{itemize}
\tightlist
\item
  a small ordered half-spectrum;
\item
  at least one output mode with multiple parent pairs;
\item
  at least two triads sharing a parent mode;
\item
  full-rank \(2\times2\) target covariances;
\item
  interaction tensors chosen first, followed by analytically generated
  realizable covariance and bispectrum targets.
\end{itemize}

After this baseline passes all analytical and Monte Carlo checks, add
rank-deficient covariances, repeated-parent triads, general
transverse-basis gauges, and exceptional modes.

\end{document}
